\section{What window certificates reach}\label{zh:sec:closer}

A window of half length $a$ resolves the zeros up to the horizon $2\pi e^{2a}$ (a heuristic, Section~\ref{zh:sec:horizon}, after Section 9 of \cite{Zhu26}), and whether the verification of the zeros up to a height $H$ implies positivity on the window of half length $\frac12\log(H/2\pi)-C$ is Problem~\ref{zh:op:horizon}; for large displacements it may fail.

\section{Open problems}\label{zh:sec:open}

\begin{problem}[The horizon as a theorem]\label{zh:op:horizon}
Determine the largest half length $a(H,\delta_0)$ such that, if all zeros with $0<\gamma\le H$ lie on the line and every zero above $H$ has $|\beta-\frac12|\le\delta_0$, Weil's form is positive on the window $[-a,a]$ with $a=a(H,\delta_0)$. Is $a(H,\delta_0)\ge\frac12\log\frac H{2\pi}-C(\delta_0)$, with a constant $C(\delta_0)$ that depends only on $\delta_0$? Status: open. The heuristic of Section~\ref{zh:sec:horizon} suggests the bound for small $\delta_0$. A window gains up to a factor $e^{2a\delta}$ on a zero at distance $\delta$ from the line, and in Table~\ref{zh:tab:det} the largest displacement is detected below the horizon at the largest heights, so a constant $C$ that does not depend on $\delta_0$ may not exist. A proof would need a quantitative form of the lower sampling inequality of Proposition~\ref{zh:prop:dict} for the ordinates, with constants that depend only on the local density of the zeros. Compare Proposition~\ref{yo:prop:barrier}, which is proved for the tent coefficient only.
\end{problem}

\begin{problem}[Decay law from the Thorin side]\label{zh:op:decay}
Derive the exponent of Zhu's decay law, $-\log\lambda_{\min}^{\rm even}(\frac12,a)=C\,N(T_*)/\log N(T_*)\,(1+o(1))$ with $T_*=2\pi e^{2a}$ and the fitted constant $C=20.13\ldots$ of his equation (5) (his Conjecture 12.1 \cite{Zhu26} states the law with $C=2\pi^2$), from Zhu's test function $\widehat f=h\,G_T(-r^2)$ of Section~\ref{zh:sec:extremal}, optimising $T$ and $h$ jointly, and explain the constant, which Zhu finds close to $2\pi^2=19.74\ldots$, through the Landau Widom plunge. Status: open. Zhu fits the law to certified upper bounds for $a\le1.2$, tests it on $a=1.4$ to $2$, and states that the constant is fitted and not derived. The only proved decay beyond the bound $O(a^{-N})$ for every $N$ of Theorem~\ref{dc:thm:ladder}(c) of Chapter~\ref{chap:dc} is his bound $\lambda_{\min}^{\rm even}(\frac12,a)\le\exp(-ae^a)$ for large $a$, under RH (Theorem 1.3), which kills the zeros only up to $2\pi e^a$, half the exponent of $T_*$, and is weaker than the law.
\end{problem}

\begin{problem}[Certified damped windows]\label{zh:op:damped}
Extend the certification of Zhu's Theorem 1.2 to the damped forms at $\alpha\in(\frac12,1)$, where under RH the margin is at least $2\xi'(\alpha)/\xi(\alpha)$, of order $\alpha-\frac12$, in place of $\lambda_{\min}(\frac12,a)$ (Theorem~\ref{zh:thm:twoparam}), and measure how far $a$ can be pushed. Status: open. No damped analogue of Zhu's Theorem 1.1 is proved here, and below $\alpha=1$ the damped pole part is not a positive rank one term (Section~\ref{zh:sec:comb}); for $\alpha\ge1$ no certificate is needed (Corollary~\ref{zh:prop:fullsymbol}). The question for the rungs of Chapter~\ref{chap:dc} is Problem~\ref{dc:prob:cert}.
\end{problem}

\begin{problem}[Comb cancellation on average]\label{zh:op:comb}
Zhu's barrier concerns the supremum of the comb $\mathcal P_a(t)=\sum_{\log n<2a}2\Lambda(n)n^{-1/2}\cos(t\log n)$. The form sees the comb only through $\int|F|^2\mathcal P_a$ with $F$ in the Paley Wiener space $PW_a$ of Proposition~\ref{zh:prop:dict}, an average over a frequency band of width about $1/a$. Determine the supremum of $\int_{|t|\ge T}|F|^2\mathcal P_a\big/\int_{|t|\ge T}|F|^2$ over $F\in PW_a$, as a function of $T$ and $a$. Status: open. Zhu's Lemma 3.2 shows that the pointwise supremum of $\mathcal P_a$ is $A_a$, which bounds the ratio from above. The average over the whole line is smaller than $A_a$ for every $F$, because $g=f*\widetilde f$ vanishes at $\pm2a$ and damps the terms with $\log n$ near $2a$; how much smaller the restricted ratio is, uniformly in $T$, is not known, and a bound of smaller order than $A_a=(4+o(1))e^a$ would be a statement about short sums of $\Lambda(n)n^{-1/2-it}$, the object of the Fej\'er tapered sums of Proposition~\ref{mw:prop:AX}.
\end{problem}

\subsection*{Notes}

The theorems of Zhu quoted in this chapter are from \cite{Zhu26} (version 2 of 2 September 2026), which the chapter reads in the language of the book and does not reprove; they and his heuristics (the horizon and the plunge rate of his Section 9, the parity of the ground state) are credited where they are stated. Yoshida's form is \cite{Yoshida92}, Weil's criterion is \cite{Weil52,Bombieri}, and the verification of the zeros to $3\times10^{12}$ is \cite{PT21}. The computations of the chapter are its own: the table of least eigenvalues in $60$ digit arithmetic (Table~\ref{zh:tab:hp}), the optimisation of Zhu's product test function (Section~\ref{zh:sec:extremal}), the comb table (Table~\ref{zh:tab:comb}), the detection experiment (Table~\ref{zh:tab:det} and Figure~\ref{zh:fig:horizon}) and the damping numbers (Table~\ref{zh:tab:cross}). \texttt{code/\allowbreak README.md} lists the scripts.
